\subsection{OPERATORS ON A TENSOR PRODUCT; TENSOR PRODUCTS AS MULTI-MODULES}

With the hypotheses and notation of no. 1, for every endomorphism \(u\) (resp. \(v\) ) of the A-module \(E\) (resp. \(F\) ), \(u \otimes l_F\) (resp. \(l_E \otimes v\) ) is an endomorphism of the \(Z\) -module \(E \otimes_A F\) ; it follows immediately from (5) (no. 2) that the mapping \(u \mapsto u \otimes l_F\) (resp. \(v \mapsto l_E \otimes v\) ) is a ring homomorphism

\[
\operatorname{End} _ {\mathbf {A}} (\mathbf {E}) \rightarrow \operatorname{End} _ {\mathbf {Z}} (\mathbf {E} \otimes_ {\mathbf {A}} \mathbf {F})
\]

(resp. \(\operatorname{End}_{\mathbf{A}}(\mathbf{F}) \to \operatorname{End}_{\mathbf{Z}}(\mathbf{E} \otimes_{\mathbf{A}} \mathbf{F}))\) ; moreover,

\[
(u \otimes \mathrm{l} _ {\mathrm{F}}) \circ (\mathrm{l} _ {\mathrm{E}} \otimes v) = (\mathrm{l} _ {\mathrm{E}} \otimes v) \circ (u \otimes \mathrm{l} _ {\mathrm{F}}) = u \otimes v\tag{7}
\]

and therefore (§ 1, no. 14) E ⊗\_A F has a canonical left bimodule structure with respect to the rings End\_A(E) and End\_A(F).

This being so, suppose given on E a \(((\mathbf{B}_i^{\prime});\mathbf{A},(\mathbf{C}_j^{\prime}))\) -multimodule structure and on F a \((\mathbf{A},(\mathbf{B}_h^{\prime});(\mathbf{C}_k^{\prime}))\) -multimodule structure (§ 1, no. 14); it amounts to the same to say that ring homomorphisms \(\mathrm{B}_i^\prime \to \mathrm{End}_{\mathbf{A}}(\mathbf{E})\) , \(\mathrm{C}_j^{\prime 0} \to \mathrm{End}_{\mathbf{A}}(\mathbf{E})\) are given with pairwise permutable images, and ring homomorphisms \(\mathrm{B}_h^{\prime \prime} \to \mathrm{End}_{\mathbf{A}}(\mathbf{F})\) , \(\mathrm{C}_k^{\prime \prime 0} \to \mathrm{End}_{\mathbf{A}}(\mathbf{F})\) with pairwise permutable images. If these homomorphisms are composed respectively with the canonical homomorphisms \(\mathrm{End}_{\mathbf{A}}(\mathbf{E}) \to \mathrm{End}_{\mathbf{Z}}(\mathbf{E} \otimes_{\mathbf{A}} \mathbf{F})\) and \(\mathrm{End}_{\mathbf{A}}(\mathbf{F}) \to \mathrm{End}_{\mathbf{Z}}(\mathbf{E} \otimes_{\mathbf{A}} \mathbf{F})\) defined above, it is seen (taking account of (7)) that ring homomorphisms

\[
\mathrm{B} _ {i} ^ {\prime} \rightarrow \operatorname{End} _ {\mathbf {Z}} (\mathrm{E} \otimes_ {\mathrm{A}} \mathrm{F}), \quad \mathrm{C} _ {j} ^ {\prime 0} \rightarrow \operatorname{End} _ {\mathbf {Z}} (\mathrm{E} \otimes_ {\mathrm{A}} \mathrm{F})
\]

\[
\mathrm{B} _ {h} ^ {\prime \prime} \rightarrow \operatorname{End} _ {\mathbf {Z}} (\mathrm{E} \otimes_ {\mathrm{A}} \mathrm{F}), \quad \mathrm{C} _ {k} ^ {\prime \prime 0} \rightarrow \operatorname{End} _ {\mathbf {Z}} (\mathrm{E} \otimes_ {\mathrm{A}} \mathrm{F})
\]

are defined with pairwise permutable images; in other words, there has been defined on \(\mathbf{E} \otimes_{\mathbf{A}} \mathbf{F}\) a \(((\mathbf{B}_i'), (\mathbf{B}_h''); (\mathbf{C}_j'), (\mathbf{C}_k'')\) -multimodule structure; it is this multimodule which is also called the tensor product (relative to A) of the \(((\mathbf{B}_i^{\prime});\mathbf{A},(\mathbf{C}_j^{\prime}))\) -multimodule E and the \((\mathbf{A},(\mathbf{B}_h^{\prime\prime});(\mathbf{C}_k^{\prime\prime}))\) -multimodule F. This multi-module is the solution of a universal mapping problem analogous to that considered in no. 1; to be precise:

PROPOSITION 3. Let G be a \(((\mathbf{B}_i'), (\mathbf{B}_h''); (\mathbf{C}_j'), (\mathbf{C}_k''))-multimodule.\)

\begin{enumerate}
\def\labelenumi{(\alph{enumi})}
\tightlist
\item
  Let \(g\) be a linear mapping of the multimodule \(\mathbf{E} \otimes_{\mathbf{A}} \mathbf{F}\) into \(\mathbf{G}\) . The mapping \(f: (x, y) \mapsto g(x \otimes y)\) of \(\mathbf{E} \times \mathbf{F}\) into \(\mathbf{G}\) is \(\mathbf{Z}\) -bilinear and satisfies relations (1) of no. 1 and the conditions
\end{enumerate}

\[
\left\{ \begin{array}{l l} f (\mu_ {i} ^ {\prime} x, y) = \mu_ {i} ^ {\prime} f (x, y), & f (x \nu_ {j} ^ {\prime}, y) = f (x, y) \nu_ {j} ^ {\prime} \\ f (x, \mu_ {h} ^ {\prime \prime} y) = \mu_ {h} ^ {\prime \prime} f (x, y), & f (x, y \nu_ {k} ^ {\prime \prime}) = f (x, y) \nu_ {k} ^ {\prime \prime} \end{array} \right.\tag{8}
\]

for all \(x \in \mathbf{E}, y \in \mathbf{F}\) , \(\mu_i' \in \mathbf{B}_i'\) , \(\nu_j' \in \mathbf{C}_j'\) , \(\mu_h'' \in \mathbf{B}_h''\) , \(\nu_k'' \in \mathbf{C}_{k}''\) , \(i, j, h, k\) arbitrary.

\begin{enumerate}
\def\labelenumi{(\alph{enumi})}
\setcounter{enumi}{1}
\tightlist
\item
  Conversely, let \(f\) be a \(\mathbf{Z}\) -bilinear mapping of \(\mathbf{E} \times \mathbf{F}\) into \(\mathbf{G}\) satisfying conditions (1) (no. 1) and (8). Then there exists one and only one linear mapping \(g\) of the multimodule \(\mathbf{E} \otimes_{\mathbf{A}} \mathbf{F}\) into the multimodule \(\mathbf{G}\) such that \(f(x, y) = g(x \otimes y)\) for \(x \in \mathbf{E}, y \in \mathbf{F}\) .
\end{enumerate}

Assertion (a) follows immediately from the definition of the multimodule structure on \(\mathbf{E} \otimes_{\mathbf{A}} \mathbf{F}\) , since for example \((x \otimes y)\nu_j' = (x\nu_j') \otimes y\) . To prove (b), we note first that Proposition 1 of no. 1 gives the existence and uniqueness of a \(\mathbf{Z}\) -linear mapping \(g\) such that \(g(x \otimes y) = f(x, y)\) for \(x \in \mathbf{E}\) , \(y \in \mathbf{F}\) ; all that is needed is to verify that \(g\) is linear for the multimodule structures. As the elements \(x \otimes y\) generate the \(\mathbf{Z}\) -module \(\mathbf{E} \otimes_{\mathbf{A}} \mathbf{F}\) , it suffices to verify the relations \(g(\mu'(x \otimes y)) = \mu_i' g(x \otimes y)\) and their analogues; but this follows immediately from the formula \(g(x \otimes y) = f(x, y)\) and relations (8).

SCHOLIUM. An element of \(\mathbf{E} \otimes_{\mathbf{A}} \mathbf{F}\) may in general be written in several ways in the form \(\sum_{i} (x_i \otimes y_i)\) , where \(x_i \in \mathbf{E}\) and \(y_i \in \mathbf{F}\) ; but to define a linear mapping \(g\) of the multimodule \(\mathbf{E} \otimes_{\mathbf{A}} \mathbf{F}\) into a multimodule \(\mathbf{G}\) , there is no need to verify that, if \(\sum_{i} (x_i \otimes y_i) = \sum_{j} (x_j' \otimes y_j')\) , then \(\sum_{i} g(x_i \otimes y_i) = \sum_{j} g(x_j' \otimes y_j')\) ; it suffices to be given \(g(x \otimes y)\) for \(x \in \mathbf{E}\) and \(y \in \mathbf{F}\) and to verify that \((x, y) \mapsto g(x \otimes y)\) is \(\mathbf{Z}\) -bilinear and satisfies conditions (1) (no. 1) and (8).

Let \(\mathbf{E}'\) be a \(((\mathbf{B}_i'),\mathbf{A},(\mathbf{C}_j')\) -multimodule, \(\mathbf{F}'\) an \((\mathbf{A},(\mathbf{B}_h'');(\mathbf{C}_k'')\) -multimodule and \(u:\mathbf{E}\to \mathbf{E}',v:\mathbf{F}\to \mathbf{F}'\) linear mappings of multimodules; it follows immediately from the definitions (no. 2) that \(u\otimes v\) is a linear mapping of the multimodule \(\mathbf{E}\otimes_{\mathbf{A}}\mathbf{F}\) into the multimodule \(\mathbf{E}'\otimes_{\mathbf{A}}\mathbf{F}'\) .

With E always denoting a right A-module, let \(_{s}A_{d}\) denote the ring A considered as an (A, A)-bimodule ( \(\S\) 1, no. 14, Example 1); by the above, the tensor product \(\mathbf{E} \otimes_{\mathbf{A}} (\mathbf{s}\mathbf{A}_{d})\) has a canonical right A-module structure such that \((x \otimes \lambda)\mu = x \otimes (\lambda\mu)\) for \(x \in E\) , \(\lambda \in A\) , \(\mu \in A\) . The mapping \((x, \lambda) \mapsto x\lambda\) of \(\mathbf{E} \times (\mathbf{s}\mathbf{A}_{d})\) into E is Z-bilinear and satisfies conditions (1) (no. 1) and (8)

(where, in the latter, the \(\mathbf{B}_i'\) , \(\mathbf{C}_j'\) and \(\mathbf{B}_h''\) are absent and the family \((\mathbf{C}_k'')\) reduces to \(\mathbf{A}\) ); hence (Proposition 3), there exists an A-linear mapping \(g\) (called canonical) of \(\mathbf{E} \otimes_{\mathbf{A}} (s\mathbf{A}_d)\) into \(\mathbf{E}\) such that \(g(x \otimes \lambda) = x\lambda\) for \(x \in \mathbf{E}\) , \(\lambda \in \mathbf{A}\) .

PROPOSITION 4. If E is a right A-module, the mapping \(h: x \mapsto x \otimes 1\) of E into \(\mathbf{E} \otimes_{\mathbf{A}} (s\mathbf{A}_d)\) is a right A-module isomorphism, whose inverse isomorphism \(g\) is such that \(g(x \otimes \lambda) = x\lambda\) for \(x \in \mathbf{E}, \lambda \in \mathbf{A}\) .

If \(g\) is the canonical mapping, \(g \circ h\) is the identity mapping \(1_{\mathbf{E}}\) and \(h \circ g\) coincides with the identity mapping of \(\mathbf{E} \otimes_{\mathbf{A}} (s\mathbf{A}_d)\) onto itself for elements of the form \(x \otimes y\) , which generate the latter \(\mathbf{Z}\) -module; hence the conclusion.

We shall normally write \(E \otimes_{A} A\) instead of \(E \otimes_{A} (s A_{d})\) and shall often identify \(E \otimes_{A} A\) with E by means of the above canonical isomorphisms. Observe that, if E also has a (left or right) B-module structure which is compatible with its right A-module structure, g and h are also isomorphisms for the B-module structures on E and \(E \otimes_{A} A\) (and hence multimodule isomorphisms).

Now let F be a left A-module; \((sA_{d}) \otimes_{A} F\) (also denoted by A \(\otimes_{A} F\) ) then has a canonical left A-module structure and as in Proposition 4 a canonical isomorphism is defined of A \(\otimes_{A} F\) onto F mapping \(\lambda \otimes x\) to \(\lambda x\) , and its inverse isomorphism is \(x \mapsto 1 \otimes x\) .

In particular there exists a canonical isomorphism of the (A, A)-bimodule \((_{s}\mathbf{A}_{d})\otimes_{\mathbf{A}}(_{s}\mathbf{A}_{d})\) onto \({}_{s}\mathbf{A}_{d}\) which maps \(\lambda \otimes \mu\) to \(\lambda \mu\) .

